\appendix{Фрагменты исходного кода программы}

В приложении приведены фрагменты исходного кода основных компонентов платформы Easyappoint: серверной части, клиентского приложения, сервиса уведомлений и миграции базы данных.

\subsection*{Фрагмент класса \texttt{BookingService}}
\addcontentsline{toc}{subsection}{Фрагмент класса BookingService}

Фрагмент демонстрирует создание бронирования, проверку роли клиента, проверку пересечений, проверку доступности слота, сохранение записи и публикацию события.

\lstinputlisting[language=Java, frame=none, firstline=42, lastline=103]{../../apps/data-service/src/main/kotlin/rodi/easyappoint/dataservice/bookings/service/BookingService.kt}

\subsection*{Фрагмент класса \texttt{AvailabilitySlotCalculator}}
\addcontentsline{toc}{subsection}{Фрагмент класса AvailabilitySlotCalculator}

Фрагмент демонстрирует построение кандидатных слотов и исключение занятых временных окон.

\lstinputlisting[language=Java, frame=none, firstline=8, lastline=51]{../../apps/data-service/src/main/kotlin/rodi/easyappoint/dataservice/availability/AvailabilitySlotCalculator.kt}

\subsection*{Фрагмент контроллера \texttt{BookingController}}
\addcontentsline{toc}{subsection}{Фрагмент контроллера BookingController}

Фрагмент демонстрирует REST-маршруты для просмотра, создания, отмены и изменения бронирований.

\lstinputlisting[language=Java, frame=none, firstline=15, lastline=85]{../../apps/data-service/src/main/kotlin/rodi/easyappoint/dataservice/bookings/controller/BookingController.kt}

\subsection*{Фрагмент маршрутизации клиентского приложения}
\addcontentsline{toc}{subsection}{Фрагмент маршрутизации клиентского приложения}

Фрагмент демонстрирует ролевую маршрутизацию React-приложения для мастера, клиента и администратора.

\lstinputlisting[frame=none, firstline=27, lastline=108]{../../apps/frontend/src/App.tsx}

\subsection*{Фрагмент обработки Kafka-событий}
\addcontentsline{toc}{subsection}{Фрагмент обработки Kafka-событий}

Фрагмент демонстрирует обработку событий бронирования и Telegram-сообщений в сервисе уведомлений.

\lstinputlisting[language=Java, frame=none, firstline=10, lastline=61]{../../apps/telegram-bot/src/main/kotlin/rodi/easyappoint/telegrambot/events/TelegramBotEventProcessor.kt}

\subsection*{Фрагмент миграции базы данных}
\addcontentsline{toc}{subsection}{Фрагмент миграции базы данных}

Фрагмент демонстрирует таблицы приглашений, расписания мастера, бронирований и outbox-событий.

\lstinputlisting[language=SQL, frame=none, firstline=92, lastline=196]{../../apps/data-service/src/main/resources/db/migration/V1__init.sql}

\ifВКР{
\newpage
\addcontentsline{toc}{section}{На отдельных листах (CD-RW в прикрепленном конверте)}
\noindent
\begin{tabular}{p{5.8cm}C{4.8cm}C{4.8cm}}
   Автор ВКР & \lhrulefill{\fill} & \fillcenter\Автор \\
            \setarstrut{\footnotesize}
           & \footnotesize{(подпись, дата)} & \\
            \restorearstrut
   Руководитель ВКР & \lhrulefill{\fill} & \fillcenter\Руководитель \\
            \setarstrut{\footnotesize}
           & \footnotesize{(подпись, дата)} & \\
            \restorearstrut
   Нормоконтроль & \lhrulefill{\fill} & \fillcenter\Нормоконтроль \\
            \setarstrut{\footnotesize}
           & \footnotesize{(подпись, дата)} & \\
            \restorearstrut
\end{tabular}
\vskip 2cm
\begin{center}
\textbf{Место для диска}
\end{center}
}\fi
